\subsection{실험 대상 및 데이터셋}
본 연구에서는 오케스트레이션 구조에 따른 성능 차이를 정량적으로 분석하기 위해, 서로 다른 특성을 가진 두 가지 유형의 문서 생성 작업을 실험 대상으로 설정하였다. 첫째는 산업 현장의 전문 지식과 정형 데이터 분석 능력이 요구되는 \textit{정형 보고서 작성 작업}이며, 둘째는 문맥적 유연성과 창의적 구성 능력이 요구되는 \textit{창의적 텍스트 생성 작업}이다.

특히 정형 보고서 작성 작업의 대표 사례로 '금형-사출-사출기-단위-가동품질-통계-분석-보고서' 생성을 사용하였다. 해당 작업은 복잡한 산업 도메인의 기술적 용어와 통계적 수치를 포함하며, 엄격한 논리적 구조와 형식적 완결성이 요구된다. 이는 고정된 워크플로우를 따르는 Code 구조와 동적 경로를 탐색하는 LLM 구조, 그리고 이를 절충한 Hybrid 구조 간의 효율성 및 품질 차이를 극명하게 드러낼 수 있는 벤치마크가 된다. 각 실험은 동일한 입력 데이터셋을 기반으로 반복 수행되었으며, 구조적 차이 외의 외부 변수를 제거하여 비교의 객관성을 확보하였다.

\subsection{통제 변수 및 실험 환경}
오케스트레이션 구조라는 독립 변수가 종속 변수에 미치는 순수한 영향을 측정하기 위해 다음과 같은 통제 변수를 설정하였다.

첫째, 에이전트의 구성과 역할(Role)을 고정하였다. 모든 구조에서 설계(Planner), 집필(Writer), 검수(Reviewer), 수정(Editor)의 네 가지 에이전트가 동일하게 배치되며, 각 에이전트에게 부여된 페르소나와 책임 범위는 동일한 프롬프트 정의를 따른다. 둘째, 사용되는 LLM 모델과 기본 프롬프트 로직을 통일하였다. 특정 구조에서만 더 고성능의 모델을 사용하거나 프롬프트를 최적화함으로써 발생하는 성능 왜곡을 방지하기 위해, 모든 에이전트는 동일한 베이스 모델을 사용한다. 셋째, 입력 데이터의 형태와 최종 출력물의 요구 규격을 동일하게 설정하여, 생성된 문서의 길이($\text{chars}$)나 내용의 범위가 구조적 특성이 아닌 입력값에 의해 결정되지 않도록 통제하였다.

\subsection{평가 지표 정의}
시스템의 성능을 다각도로 분석하기 위해 자원 효율성을 측정하는 '효율성 지표'와 결과물의 완성도를 측정하는 '효과성 지표'로 나누어 정의하였다.

\begin{enumerate}
    \item \textbf{효율성 지표 (Efficiency Metrics)}
    \begin{itemize}
        \item 실행 시간 ($\text{elapsed\_sec}$): 요청 입력부터 최종 결과물 출력까지 소요된 총 시간을 측정하여 시스템의 응답성을 평가한다.
        \item 토큰 소비량 ($\text{tokens}$): 전체 프로세스에서 소모된 입력 및 출력 토큰의 총합을 측정하여 경제적 비용 및 자원 소모량을 산출한다.
    \end{itemize}
    \item \textbf{효과성 지표 (Effectiveness Metrics)}
    \begin{itemize}
        \item 품질 점수 ($\text{best\_score}$): LLM-as-Judge 방식을 통해 산출된 정량적 점수로, 생성된 문서의 논리성, 정확성, 완결성을 평가한다.
        \item 재시도 횟수 ($\text{retry\_count}$): 검수 에이전트가 반려하여 수정 에이전트에게 다시 작업을 요청한 횟수를 측정한다. 이는 오케스트레이션 구조가 목표 품질에 도달하기 위해 수행하는 내부 루프의 빈도를 의미하며, 구조의 안정성과 수렴 속도를 나타내는 지표가 된다.
    \end{itemize}
\end{enumerate}

\subsection{LLM-as-Judge 평가 방법론}
생성된 문서의 품질을 객관적으로 평가하기 위해, 인간 평가자의 주관성을 배제하고 일관된 기준을 적용할 수 있는 LLM-as-Judge 프레임워크를 도입하였다. 평가 모델은 실험에 사용된 모델보다 상위 추론 능력을 가진 모델을 선정하여 평가자로 활용하였다.

평가 프롬프트는 다음과 같은 구체적인 루브릭(Rubric)을 포함하도록 구성되었다: (1) 설계도(Plan)에 명시된 모든 필수 항목이 누락 없이 포함되었는가, (2) 산업 도메인의 기술적 용어가 정확하게 사용되었는가, (3) 문장 간의 논리적 흐름이 매끄럽고 일관성이 있는가, (4) 보고서로서의 형식적 요건을 충족하는가. 평가 모델은 각 항목에 대해 근거를 먼저 제시한 후, 최종적으로 0점에서 100점 사이의 정량적 점수를 부여하도록 설계되었다. 이러한 방식은 단순한 정성적 평가를 넘어, 구조별 평균 점수와 표준편차를 산출함으로써 통계적 유의성을 확보하는 근거가 된다.